\chapter{V/Q Scan}
\section*{What Is This Test?} A ventilation/perfusion (V/Q) scan compares air flow through the lungs with blood flow to detect mismatches.\section*{Alternative Names} VQ scan; lung scan; ventilation-perfusion scintigraphy.\section*{Principle} Inhaled and injected radioactive tracers show ventilation and perfusion on nuclear images.\section*{Why This Test Is Ordered} It evaluates suspected pulmonary embolism and selected chronic pulmonary vascular or ventilatory disorders, especially when CT angiography is unsuitable.\section*{Primary vs Secondary Test} It is a secondary imaging study selected by clinical probability, chest findings, kidney function, and pregnancy.\section*{How This Test Is Performed} The patient inhales a tracer and receives an intravenous tracer, followed by images of the lungs.\section*{What Preparation Is Needed} Usually little preparation. Report pregnancy, breastfeeding, kidney disease, and inability to lie flat.\section*{Understanding Results} Normal, low, intermediate, or high-probability patterns are interpreted with clinical risk; nondiagnostic results need another plan.\section*{Follow-up Tests} CT pulmonary angiography, leg ultrasound, ECG, blood tests, or specialist assessment may follow.\section*{References} \begin{enumerate}\item MedlinePlus. V/Q Scan.\item Society of Nuclear Medicine and Molecular Imaging. V/Q procedure standard.\end{enumerate}\clearpage\section*{Understanding Results and Follow-up}
The laboratory report should identify the specimen, method, units, reference interval or decision limit, and whether the result is positive, negative, abnormal, or inconclusive. Reference intervals vary with age, sex, pregnancy, specimen type, assay, and laboratory; a result outside a range is not automatically a diagnosis, and a result within a range does not always exclude disease. Discuss unexpected results with the ordering clinician, who can decide whether repeat or confirmatory testing, treatment, or referral is appropriate. Do not change medicines or delay urgent care based on an isolated result.
\section*{Regulatory Status and Authorization Date}This chapter describes a test category, not a specific commercial product. FDA, CDSCO, EU IVDR, and MHRA approval or registration dates therefore cannot be assigned without the assay manufacturer, product name, instrument, and intended use. For laboratory-developed tests, an individual agency approval date may not exist. See the Preface for the interpretation of regulatory status.
\clearpage
